\section{Vacuum Uniqueness via Cluster Decomposition}
\label{sec:vacuum_uniqueness}

We must prove the vacuum $\Omega$ is unique. If the vacuum were degenerate, the theory would decompose into superselection sectors.

\subsection{The Cluster Property}
We have already proven the mass gap $\Delta > 0$ (via Giles-Teper and Adjoint Interpolation). The spectral gap implies exponential decay of connected correlation functions:
\begin{equation}
|\langle A(x) B(y) \rangle - \langle A \rangle \langle B \rangle| \le C e^{-\Delta |x-y|}
\end{equation}

\subsection{Equivalence to Uniqueness}
By the Haag-Ruelle scattering theory, the cluster property holds if and only if the vacuum is unique.

\begin{theorem}[Uniqueness of Vacuum]
\label{thm:unique-gibbs}
The infinite-volume vacuum state $|\Omega\rangle$ is unique (non-degenerate).
\end{theorem}

\begin{itemize}
\item If there were multiple vacua $\Omega_i$, correlations would approach a constant non-zero value at infinity (long-range order).
\item Since we proved exponential decay (finite correlation length $\xi = 1/\Delta < \infty$), the vacuum must be unique.
\end{itemize}
